% Ad Atticum 13.45 (ad-atticum-13-45)
% Source: https://raw.githubusercontent.com/PerseusDL/canonical-latinLit/master/data/phi0474/phi057/phi0474.phi057.perseus-lat1.xml
% Fetched by scripts/fetch_latin.py

% Dateline (Perseus): Scr. in Tusculana iii Id. Sext. a. 709 (45).

\ciceroLetterOpener{CICERO ATTICO salutem}

\ciceroSection{1} fuit apud me Lamia post discessum tuum epistulamque ad me attulit missam sibi a Caesare. quae quamquam ante data erat quam illae Diocharinae, tamen plane declarabat illum ante ludos Romanos esse venturum. in qua extrema scriptum erat ut ad ludos omnia pararet neve committeret ut frustra ipse properasset. prorsus ex his litteris non videbatur esse dubium quin ante eam diem venturus esset, itemque Balbo cum eam epistulam legisset videri Lamia dicebat. dies feriarum mihi additos video sed quam multos fac, si me amas, sciam. de Baebio poteris et de altero vicino Egnatio.

\ciceroSection{2} quod me hortaris ut eos dies consumam in philosophia explicanda, currentem tu quidem; sed cum Dolabella vivendum esse istis diebus vides. quod nisi me Torquati causa teneret satis erat dierum ut Puteolos excurrere possem et ad tempus redire.

\ciceroSection{3} Lamia quidem a Balbo, ut videbatur, audiverat multos nummos domi esse numeratos quos oporteret quam primum dividi, magnum pondus argenti; auctionem XIII. xlv M. TVLLI CICERONIS praeter praedia primo quoque tempore fieri oportere. scribas ad me velim quid tibi placeat. equidem si ex omnibus esset eligendum, nec diligentiorem nec officiosiorem nec me hercule nostri studiosiorem facile delegissem Vestorio, ad quem accuratissimas litteras dedi; quod idem te fecisse arbitror. mihi quidem hoc satis videtur. tu quid dicis? unum enim pungit ne neglegentiores esse videamur. exspectabo igitur tuas litteras.
